% Chapter 2 -- translation by Doug Burkett & Tim Hutcheson (ca. 2001),
% lightly copy-edited; a few printed values corrected against the
% Spanish original (pr2, pr3, is4 in Problems 008 and 004).

\chapter{Numerical simulation of direct current (DC) circuits}

\section{Before starting a simulation}

Since Symbulator gives you the simulation results in variables
stored in the current folder, it is good practice to delete
unnecessary variables in that folder before simulating. This
prevents confusing both you and Symbulator with answers of a
previous simulation. Although it is not necessary to do so, this
practice will ensure better simulation results.

\subsection{Current folder}

The current folder is the folder to which the calculator is
currently set. Usually, the current folder is the main folder of the
calculator, called MAIN. You can use any folder you wish, but I
recommend that you use MAIN unless you have a specific reason to use
a different folder.

The current folder is shown on the lower left screen of the
calculator in small text.

\subsection{Unnecessary variables}

Unnecessary variables are the variables in the current folder which
are not needed to run the simulation. One example of unnecessary
variables are the results of a previous simulation.

\subsection{Procedure}

If you are familiar with the use of your calculator, you already
know how to delete variables. A complete explanation can be found in
the User's Guide. In any case, use these steps to delete the
variables:

Identify the current folder. The current folder name is shown in the
lower left corner of the calculator screen. Obviously, the
calculator must be on. If you do not see the folder name, press
\key{HOME}, which displays the Home screen.

The home screen has several sections. The upper section is called
the Toolbar. The central section is called the History area. Entries
and results are shown here. The line below the history area is
called the entry line, because this is where you enter calculator
commands. The last section, below the entry line, is called the
Status line. The Status line shows the state of several calculator
modes, including the current folder. MAIN is shown if the main
folder is current.

Start Var-Link. The Var-Link screen shows the folders and variables
which you have on your calculator. The Var-Link screen is used to
manipulate variables, which includes deleting them. To display the
Var-Link screen on the TI-89, press \key{2nd} \key{$-$}; note that
VAR-LINK is shown above the \key{$-$} key.

Open the current folder. The Var-Link screen first shows the folders
on your calculator. At the upper right of the calculator keyboard
are four blue arrow keys, called the cursor keys. Press the blue
\key{UP} and \key{DOWN} cursor keys to highlight the current folder
name. If an arrow is shown to the right of the folder name, that
folder contains variables. If no arrow is shown, there are no
variables in the folder. If there are no variables, you need do
nothing more in Var-Link: press \key{HOME} to return to the home
screen. Otherwise, open the folder by pressing the blue \key{RIGHT}
key.

Delete the unnecessary variables. To delete the unnecessary
variables, first mark them all by using the blue \key{UP} and
\key{DOWN} keys to highlight the variable names, then press \key{F4}
to mark them. A small check-mark appears to the left of the variable
name. When all the unnecessary variables are marked, press the
\key{BACKSPACE} key to delete them. The \key{BACKSPACE} key is above
the \key{T} key on the TI-89, and is marked with a left-pointing
arrow. A dialog box is shown which asks if you want to delete the
variables. Press \key{ENTER} to delete them.

To delete all variables from the MAIN folder, select the folder,
then press \key{BACKSPACE}. All the variables in MAIN will be
deleted, but the folder will remain. The MAIN folder cannot be
deleted because it is the `main' calculator folder. An error message
is displayed, indicating that the folder cannot be deleted, but the
variables are deleted, nonetheless. This method is recommended,
because it is fast and effective. However, if you use this method,
you must avoid storing variables in MAIN which you want to keep.

\section{First step for all simulations}

The Symbulator offers four different circuit simulation analyses.
Each analysis uses a different `door' and operates with a different
logic. Nevertheless, all of them have something in common: they
require that the student knows how to correctly describe the
elements and nodes of the circuit. Thus, the first step for all
simulations is:

\begin{quote}
\emph{Name all the nodes and elements of the circuit.}
\end{quote}

In Chapter 1 we explained the way to name the nodes. The way to name
the elements will be explained in different chapters, as they are
used. All these names are unique, and once they are assigned, you
must be consistent in using them.

\section{Door for direct current: dc}

As mentioned, the door is the name of a program which calls other
programs to actually perform the simulation. In order to run a
direct current (DC) circuit simulation, use the door
\code{sq\char`\\dc}. The letters \code{dc} are the name of the door.
The letters \code{sq} tell the calculator in which folder the
program \code{dc} is located. Thus, entering \code{sq\char`\\dc}
tells the calculator ``execute the program \code{dc} in folder
\code{sq}''.

\section{Input data: the circuit}

When using the \code{dc} door or any of the other four doors, you
must provide the program the circuit you want to simulate, as input
data. The circuit description must be placed just after the name of
the door, between the parentheses, like this:
\code{sq\char`\\dc(circuit)}.

There are two ways to do this. The first way is to put the text
which describes the circuit between the parentheses, for example:

\begin{entry}
sq\dc("description of the circuit")
\end{entry}

The second way is to place a variable name, in which the text has
been stored, between the parentheses, for example:

\begin{entry}
"description of the circuit" -> variable
sq\dc(variable)
\end{entry}

Both methods work well. The second method has the advantage that the
circuit is saved in memory as a variable, so you may use it again
later, or change it to run a slightly different simulation.

\section{Answers}

\subsection{Philosophy}

The philosophy of the Symbulator is to give you all the simulation
results. These results are not displayed automatically, but are
stored in variables with intuitive names. The idea behind this is
that you need only display the results of interest.

In the case of DC circuit simulation, the results are the voltages,
the currents and the powers. Let us see:

\subsection{Node voltages}

The Symbulator calculates node voltages with respect to the
reference node, in volts (V). These voltages are stored in variables
whose name is \code{v} plus the name of the node. For example, if
the node is named \code{q}, the name of the variable in which the
voltage is stored is \code{vq}.

\subsection{Element voltage drops}

The Symbulator calculates the voltage drops across all the circuit
elements, defined as the voltage drop from the first node in the
description with respect to the second node, in volts (V). These
voltage drops are stored in variables whose name is \code{v} plus
the name of the element. For example, if the element is called
\code{r5}, the name of the variable in which the voltage drop is
stored is \code{vr5}.

\subsection{Currents}

The Symbulator calculates the currents through all the elements,
defined as the current from the first node in the description to the
second node, in amperes (A). The currents are stored in variables
whose name is \code{i} plus the name of the element. For example, if
the element is called \code{r5}, the name of the variable in which
the current through the element is stored is \code{ir5}.

\subsection{Powers}

The Symbulator calculates the real power dissipated by all the
elements, in watts (W). The powers are stored in variables whose
name is \code{p} plus the name of the element. For example, if the
element is called \code{r5}, the name of the variable in which the
real power consumed by the element is stored is \code{pr5}. Note
that the power is defined as if it were dissipated. In the case of
an element that generates power instead of consuming it---as is the
case of most of the sources---the power in the answer will be
negative.

\section{Short circuit}

Ideally, a short circuit is a perfect conductor which transmits any
current without voltage drop or power consumption.

\subsection{Notation}

From the point of view of the simulation, the short circuit is the
most simple element. Nevertheless, it is the least-used element. We
begin with the short circuit, because it is the easiest to describe.

\emph{What class of element is it?} It is a short circuit. To the
Symbulator, we indicate the class of element by means of the first
letter of the name that we give the element. The letter that
identifies the short circuits is \code{s}. That is to say, if the
first letter of the name of the element is \code{s}, then Symbulator
knows that the element is a short circuit.

\emph{How is the element named?} Since a circuit can have several
elements of the same class, we must give a unique name to each
element, so that Symbulator can distinguish the different elements.
Therefore, if a circuit has two short circuits, we must give each
one a different name. However, both names must begin with \code{s},
because both represent short circuits. For example, we can name one
of them \code{s1} and another \code{s2}. Another alternative would
be to name one \code{sa} and another \code{sb}. Any name is valid
which begins with \code{s} and that is made up of numbers and/or
letters. Once the user has defined the name of an element, the
Symbulator will use that name for all answers related to that
element.

\emph{Where is the element connected?} A short circuit has two ends,
each one of them connected to a different node. You must specify
both nodes to which the short circuit is connected. In the
Symbulator, the order in which the nodes are listed always indicates
the direction of current flow through the element. In the case of a
short circuit, the current is defined as flowing through the element
from the first node to the second node of the description.

I have mentioned that each circuit element is described as a row of
terms in the circuit description. Each row includes the minimum
information needed to specify the element and its connections. So, a
short circuit is defined as:

\begin{entry}
sname, first node, second node
\end{entry}

This row condenses, in simple form, the minimum information
necessary to describe the short circuit. Note that there are only
three elements in the row: a name that begins with \code{s}, the
first node and the second node. There is a fourth specification in
the row, which is not evident at first sight: the order in which the
nodes appear specifies the direction in which the user wishes to
define current flow through the element. Implicit in the description
is the order: ``define the short circuit current through element
\code{sname} as flowing from the first node to the second node.''

\subsection{Symmetry of the description}

Once Symbulator starts, the circuit description in text form is
converted into a matrix. Like any other computer, your calculator
does not allow asymmetric matrices: all the rows must have the same
number of elements (dimension). If one row has $n$ terms, then all
rows must have $n$ terms, otherwise an error occurs.

I have mentioned that three terms are required to describe a short
circuit. Other circuit elements such as sources require four terms.
Still others, such as energy storage elements, require five terms.

As all the rows must have the same number of terms, you need to
identify which row has the most terms. The shorter rows must be
filled with zeros so that they are the same length as the longest
row.

For example, in a circuit description with a longest row of four
terms, you need to add a zero to the row which describes a short
circuit, so that it has four terms, also:

\begin{entry}
sname, first node, second node, 0
\end{entry}

As another example, with a circuit description whose longest row has
five terms, you must add two zeros to the short circuit description
row:

\begin{entry}
sname, first node, second node, 0, 0
\end{entry}

In the case of other circuit elements, it will sometimes be
necessary to add zeros at the end of the rows to maintain the
symmetry of the internal matrix.

\subsection{Related answers}

I have said that each node voltage is stored in a variable at the
end of the simulation. In addition, each circuit element will
generate results related to that element, and they will be stored in
variables with intuitive names.

In the case of a short circuit, a single related result is given:
the current through the short circuit. Since the voltage drop is
zero, no power is dissipated, so there are no other related results.

\section{Resistors}

Ideally, a resistor is an imperfect conductor which resists to
greater or lesser degree the flow of current. The resistor causes a
voltage drop and therefore dissipates power.

\subsection{Notation}

The information necessary to completely describe a resistor is as
follows:

\emph{What is the element type?} It is a resistor, so it is
identified with the letter \code{r}.

\emph{How is the resistor named?} In order to distinguish a given
resistor from any other, each resistor is given a unique name. The
name must begin with \code{r}.

\emph{Where is the resistor connected?} A resistor has two
terminals, each connected to a different node. You must tell
Symbulator the name of the nodes to which the terminals are
connected, in the order in which you want the current flow and
voltage drop defined. The current is defined as flowing through the
resistor from the first node to the second node. The voltage drop is
defined as the voltage of the first node with respect to the second
node.

\emph{What is the resistor value?} The value must be given to the
Symbulator in ohms. For a DC analysis, the value may be a real
number or a variable.

So, resistors are defined as:

\begin{entry}
rname, node 1, node 2, value in ohms
\end{entry}

Note that the row which describes the resistor has four terms. If
another row of the description matrix has five terms, then you must
add a zero to the end of the resistor description, for example:

\begin{entry}
rname, node 1, node 2, value, 0
\end{entry}

The resistor value must be entered in ohms; for example, if the
resistance is 1~k$\Omega$, you must enter the value as 1000 ohms, or
the results will be wrong.

For simulations other than DC analysis, the resistance value can be
a phasor, or a function of frequency or time.

\subsection{Conductance}

A conductance and a resistance are the same element analyzed from
two different perspectives. The resistance is analyzed from the
point of view of how much it resists current flow, and is measured
in ohms. Conductance is analyzed from the point of view of how much
it allows current flow, and is measured in siemens or mhos.

If we have a value of resistance and we want to convert it to
conductance, we find the reciprocal of the resistance. To find the
reciprocal, raise the value to the $-1$ power or find one divided by
the value. For example, a resistance of 5 ohms is equivalent to a
conductance of $5^{-1}$ or $1/5 = 0.2$ siemens.

For the purposes of simulation, conductances must be specified as
resistances. The value must be entered in ohms, not siemens. You can
perform the conversion directly in the resistor description:

\begin{entry}
rname, node 1, node 2, 1/(value in siemens)
\end{entry}

or:

\begin{entry}
rname, node 1, node 2, (value in siemens)^-1
\end{entry}

\subsection{Related results}

In the case of a resistor, three related results are given: the
current through the resistance, the voltage drop and the real power
dissipated by the resistance. The current is stored in variable
\code{iname}, where \code{name} is the resistor name. The voltage
drop is stored in a variable called \code{vname}. The dissipated
power is stored in a variable called \code{pname}.

\section{Current source}

An ideal current source is an element which forces a defined current
flow through itself, and generates or consumes power.

\subsection{Notation}

The information necessary to define a current source is:

\emph{What is the element type?} It is a current source. Current
sources are defined with the letter \code{j}.

\emph{What is the name of the source?} To distinguish a particular
current source from other sources in the circuit, it must be given a
unique name. The name begins with \code{j}.

\emph{Where is the source connected?} The current source has two
terminals, each of which is connected to a different circuit node.
You must tell Symbulator the nodes to which the current source is
connected. The order in which you list the nodes specifies the
defined direction of current flow: the current flows from the first
node to the second node. In the circuit schematic, the current flow
direction is shown by the arrow in the source symbol. The order of
the nodes also implicitly defines the sign of the voltage drop
across the source: the voltage drop is the voltage of the first node
with respect to the second node.

\emph{What is the value of the current source?} The value of the
current source is specified in amperes (A). In a DC analysis, the
value of an independent current source is a real number or a
variable. The value of a dependent current source is specified as an
expression. The expression may be a number and a variable name. The
variable is one of the future simulation results. The examples will
show this case more clearly.

So, a current source is defined as:

\begin{entry}
jname, node 1, node 2, value in amperes
\end{entry}

As mentioned, the direction of the current flow is defined as node 1
to node 2. For example, if the value of the current source is 10,
then 10~A flow through the source from node 1 to node 2.

Note that the row that defines a current source has four terms. You
need to add a zero to the end of the definition if at least one
other row has five terms. For example:

\begin{entry}
jname, node 1, node 2, value, 0
\end{entry}

Make sure that you enter the value in A. For example, if the current
source is specified as 10~mA, you must enter a value of 0.01~A. If
you enter the value in mA, the results will be wrong.

For simulations other than DC analysis, the current source value may
be a phasor, or a function of time or frequency.

The fact that dependent sources are defined with the same notation
as independent sources is one of the great advantages of Symbulator
compared to other simulators, which require different names and
special techniques using control nodes or branches to define
dependent sources.

\subsection{Related results}

In the case of a current source, three related results are given:
the current through the source, the voltage drop across the source
and the real power consumed by the source. The current is stored in
a variable called \code{iname}, where \code{name} is the current
source name. The voltage drop is stored in a variable called
\code{vname}. The consumed real power is stored in a variable called
\code{pname}.

It may seem pointless that the value of a current source is given as
a result. It could be said that the value must be known before the
simulation, so giving it as a result does not contribute any new
information. This is true with numerical simulators. However, it is
different in the case of Symbulator, which offers 100\% symbolic
capabilities, because the value of the current source may be unknown
at the start of the simulation: the value may be a variable. For
that reason, the current would not be known, which makes it a very
important result for the user. That is the reason why Symbulator
stores the values of all results in variables.

Depending on the circuit configuration, a current source may consume
power. However, most often current sources generate power. I
mentioned above that power results are positive for dissipated
power. If an element, including a current source, generates power,
then the actual generated power is the negative of the value stored
in the result variable. It is important to keep this in mind when
you interpret the answers.

For example, if the answer in the variable \code{pname} has a value
of $-10$, this means that the source dissipated $-10$ watts. This is
the same as saying that the source generated 10 watts. As another
example, if \code{pname} has a value of 15, this means that the
source dissipated 15 watts, which is the same as saying that the
source generated $-15$ watts.

\problem{001}

Now I will solve the first problem. Since this is the first, the
procedure will be explained in detail. In following problems, less
explanation will be necessary, and it will be omitted gradually.

\emph{Problem: Given the following circuit, obtain the voltages of
all the nodes.}

Solution:

The first step, as always, consists of naming all the circuit nodes
and elements. This circuit has three nodes and five elements. This
is always the first step of all simulations, regardless of the type
of analysis.

One of the three nodes must be specified as the reference node. The
usual practice for selecting the reference node is to choose that
node which has the most sources connected to it. Actually, any node
can be the reference node. For this problem we choose the bottom
node as the reference, and we give it the name 0. We name the other
two nodes 1 (the left node) and 2 (the right node).

We name the elements as follows. The 2 ohm resistance is \code{r1},
the 5 ohm resistance is \code{r2}, the 1 ohm resistance is
\code{r3}, the 3.1~A current source is \code{j1} and the $-1.4$~A
current source is \code{j2}. The nodes of each current source must
be entered in the description in an order such that they indicate
the direction of current flow shown in the drawing. The nodes of the
resistors can be entered in any order.

The first step is done.

The second step is to write the circuit description text and to
start the simulation. This step could be divided in two steps: first
write the description and store it in a variable, then run the
simulation with that variable. However, I prefer to perform these
two steps as one.

The simulation is started by typing, in the calculator entry line,
the command \code{sq\char`\\dc} followed by the circuit description
between parentheses. The calculator interprets this entry as:
``Execute the program \code{dc} in folder \code{sq}, using the
circuit description I am giving you between parentheses''.

First we ensure that the calculator is ready for a new entry by
pressing the \key{CLEAR} key. Next we type this in the entry line:

\begin{entry}
sq\dc("j1,0,1,3.1; r1,1,0,2; r2,1,2,5; r3,2,0,1; j2,2,0,-1.4")
\end{entry}

Enter the numbers by pressing the number keys. Enter the
\code{\char`\\} by pressing \key{2nd} \key{2}. Enter the letters by
pressing the \key{ALPHA} key followed by the required letter key.
The letter is shown in purple above the key.

Note that the `$-$' sign in the value of the \code{j2} current
source is the negation symbol, which is entered by pressing
\key{($-$)}, not \key{$-$}.

It is also important that you understand what the negative current
source value means. In the circuit description, we have specified
the flow of current from node 2 to node 0. Since the value is
negative, this means that a current of 1.4~A flows from node 0 to
node 2.

Since no element in this description has five terms, we need not add
zeros to any of the description rows.

When the description is entered, press the \key{ENTER} key to start
the simulation. As the simulation runs, various status messages are
displayed to indicate the simulation progress. After a few seconds,
the word `Done' is shown in the history display to indicate that the
simulation is complete. The actual time to complete the simulation
depends on the particular calculator model and free memory, but it
took about 13 seconds on my calculator.

The simulation results are stored in variables in the current
folder. You can use Var-Link to display the results. If you have
followed the procedure to delete unnecessary variables, you will see
this in the Var-Link display:

\begin{itemize}
\item The currents through all the elements, in amperes: \code{ij1},
  \code{ij2}, \code{ir1}, \code{ir2}, \code{ir3}. These currents are
  defined as flowing from the first node to the second node of each
  element description.
\item The voltage drops for all elements, in volts: \code{vj1},
  \code{vj2}, \code{vr1}, \code{vr2}, \code{vr3}. The voltage drops
  are defined as the first described node voltage with respect to
  the second node.
\item The real powers consumed by all the elements, in watts:
  \code{pj1}, \code{pj2}, \code{pr1}, \code{pr2}, \code{pr3}.
\item The node voltages with respect to the reference node, in
  volts: \code{v0}, \code{v1} and \code{v2}.
\end{itemize}

If the user had not erased the unnecessary variables, they will also
be there when the user views the folder using Var-Link.

Return to the entry line by pressing the \key{HOME} key, then press
\key{CLEAR} to clear the entry line.

The problem asked for all the node voltages. To display these
voltages, we enter each node voltage variable name in the entry line
and press \key{ENTER}, and the voltage is shown in the history
display.

If we enter \code{v0} in the entry line and press \key{ENTER}, 0 is
shown in the history display. This is the expected result, since the
reference node voltage must always be zero.

If we enter \code{v1} in the entry line and press \key{ENTER}, 5. is
shown, indicating that node 1 is at 5 volts.

If we enter \code{v2} in the entry line and press \key{ENTER}, 2. is
shown, indicating that node 2 is at 2 volts.

We have finished the problem, since we have found the values of the
requested unknowns.

\problem{002}

\emph{Problem: Given the following circuit, find the voltages of all
the nodes.}

Solution:

The only new concept in this circuit compared to the previous
problem is that the resistors are specified as conductances instead
of resistances. As already explained above, we can enter the
conductances directly in the circuit description.

As always, the first step consists of naming all the nodes and
elements of the circuit. In this circuit there are four nodes, and
we need to choose one of them as the reference node. We chose the
bottom node, because it has the most sources connected to it. The
circuit has eight elements: three independent current sources and
five conductances. We will simulate the conductances as resistors.
The node and element names will be shown in the next step.

Note that the actual node and element names are a matter of personal
preference.

The second step is to enter the simulation command and circuit
description in the calculator entry line, where the circuit
description is entered, as always, between the parentheses of the
simulation command:

\begin{entry}
sq\dc("j1,0,1,-8; j2,2,1,-3; j3,3,0,-25; r1,1,2,1/3;
      r2,1,3,1/4; r3,2,3,1/2; r4,2,0,1; r5,3,0,1/5")
\end{entry}

Remember that the negative signs in front of the current source
values are entered with the negation key \key{($-$)}, not the
subtraction key \key{$-$}.

After entering the circuit description, press \key{ENTER} to run the
simulation. The calculator screen shows the status messages as the
simulation runs, then the word `Done' is shown in the history
display when the simulation is finished. This simulation takes about
20 seconds.

The simulation results are stored in the current folder with the
appropriate names. We do not need to display the reference node
voltage, because it is zero.

Enter \code{v1} in the entry line and press \key{ENTER}, and the
voltage of node 1 is shown: 1.

Enter \code{v2} in the entry line and press \key{ENTER}, and the
node 2 voltage is shown: 2.

Enter \code{v3} in the entry line and press \key{ENTER}, and the
node 3 voltage is shown: 3.

These are the correct answers. Before trying more problems, we will
learn some new concepts.

\section{Modes of operation}

The TI-89, in my opinion, is the best calculator that has ever been
made. One excellent aspect of the calculator is the intelligent way
in which it handles exact and approximate numbers.

If we press the \key{MODE} key, the calculator settings are
displayed. The settings are divided in three groups, which can be
seen by pressing \key{F1}, \key{F2} and \key{F3}. Press \key{F2} to
see the second group. Near the bottom of the screen you will see the
Exact/Approx setting, which is Auto. This setting controls the way
in which the calculator handles exact and approximate numbers. Press
\key{HOME} to return to the home screen, and press \key{CLEAR} to
clear the entry line. Now we will see how this mode setting affects
Symbulator operation.

The Symbulator handles the Exact/Approx setting automatically,
setting it to Auto. When your calculator is set to Auto, the number
2 is interpreted as an exact value, while 2.\ (with a decimal point)
is interpreted as an approximate value. The calculator considers as
approximate any number with a decimal point. Let us see this in two
examples.

In the entry line, type \code{2/3} and press \key{ENTER}. The result
is shown in the history area as \code{2/3}, which is the exact
result. The calculator has interpreted the numbers as exact, and
does not approximate the result.

Now, type \code{2./3} and press \key{ENTER}. The result shown in the
history display is \code{.666666667}. This answer, which would raise
the hair of any reader of the Apocalypse, means that the calculator
has interpreted the fraction as composed of an approximate number
and an exact number. Therefore, at least one of the numbers is
approximate, and the calculator automatically returns an approximate
decimal result.

The observant reader will note that the voltages in Problem 001
included a decimal point, while the voltages in Problem 002 did not
have the decimal point. Why? The explanation is simple. In the first
problem, the element values included approximate numbers with
decimal points (for example, 3.1 or $-1.4$), and therefore the
answers were approximate. In the second problem, all the values were
exact numbers, so the results were also exact.

This teaches us something: when the circuit values are entered as
exact numbers, the results will also be exact. And whenever the
values are entered as approximate numbers, the answers will be
approximate.

One way to turn an exact number, for example 31/10, into an
approximate value, is to use the \code{approx()} function. So, if we
enter \code{approx(31/10)}, the calculator returns the approximate
result, which is 3.1.

We can also turn an approximate number into an exact number with the
\code{exact()} calculator function. For example, \code{exact(3.1)}
returns 31/10.

\section{Dependent sources}

In simple terms, the value of a dependent source depends on another
variable in the circuit. The value of a current-dependent current
source is a function of a control current, elsewhere in the circuit.
On the other hand, the value of a voltage-dependent voltage source
is a function of a control voltage, elsewhere in the circuit.

I have said that one of the great advantages of Symbulator is that
it allows natural simulation of dependent sources, without
additional structures or nomenclature. This results from the method
that Symbulator uses to generate equations, as opposed to
conventional simulators which use methods such as modified nodal
analysis or the gyrator transform. Simulators based on these methods
require special nomenclature to distinguish independent sources from
dependent sources.

Two precautions are necessary when simulating circuits with
dependent sources. The first is to ensure that the control variable
is expressed in terms of simulation results. The second precaution
is to ensure that the current folder has no unnecessary variables,
particularly those whose names are identical to the names of results
of the current simulation that will be used to control the dependent
source.

\subsection{Voltage-dependent voltage sources}

The control voltage of a voltage-dependent voltage source can be a
node voltage, which is the voltage of the node measured with respect
to the reference node. For example, the control voltage might be
specified as \code{v1}, where \code{v1} is the voltage of node 1.
The reference voltage is always zero volts.

The control voltage may also be specified as the voltage drop across
an element. These voltage drops are included in the Symbulator
results, so a control voltage might be specified as \code{vr1},
which is the voltage drop across \code{r1}.

The control voltage may also be specified as the voltage between any
two nodes. In this case the control voltage is expressed as a
differential voltage, for example, \code{(v1-v2)}.

Let us see a problem as an example.

\problem{003}

\emph{Problem: For the circuit with two nodes shown in the figure,
find the value of the currents $i_A$, $i_B$ and $i_C$.}

Solution:

As usual, the first step is to name the circuit nodes and elements.
The circuit has two nodes. We will use the lower node as the
reference. This circuit has a voltage-dependent current source. The
control voltage is the voltage drop across the 5.6~A current source,
and is labeled $v_x$. The current through this source flows from the
lower node to the upper node. This means that when we describe it,
its voltage drop will be defined from the lower node with respect to
the upper node, a polarity which agrees with the drop $v_x$ shown in
the drawing. If we name this source \code{jx}, then the control
voltage is \code{vjx}.

The second step is to enter the simulation command and circuit
description in the calculator entry line.

The circuit specifies the direction of current flow in the sources.
But how shall we define current flow through the resistors? This is
not specified, so we must decide.

The problem asks us to find $i_A$, $i_B$ and $i_C$, which flow from
the top to the bottom of the figure. If we define the resistor
currents in the same direction, the simulation results will agree
with the figure. This means we will not have to change the signs of
the results to answer the problem questions. However, this decision
is optional, and either direction may be used. This decision just
makes it easier to interpret the simulation results.

In order to make the answers even easier to understand, we will name
the elements in the branches of $i_A$, $i_B$ and $i_C$ as \code{rA},
\code{jB} and \code{rCC}, respectively. This is optional; you may
use any names you choose. However, choosing appropriate element
names makes the results easier to interpret. Note that we cannot use
the variable name \code{rC}, because that is a reserved system name,
as pointed out in section 1.9.4.

This is the command and circuit description:

\begin{entry}
sq\dc("jx,0,1,5.6; ra,1,0,18; jb,1,0,.1*vjx; rcc,1,0,9; j2,1,0,2")
\end{entry}

The value of the dependent source could also have been described as
\code{-.1*v1}. Press the \key{ENTER} key to start the simulation.

After pressing \key{ENTER}, the simulation status messages are
shown, and soon `Done' is shown in the history display, indicating
that the simulation is complete. This simulation took 12 seconds on
my calculator. The simulation results for the element voltages,
currents and powers are stored in the current folder.

Enter \code{ira} in the entry line and press \key{ENTER}, then 3.\
is shown in the history display, which is the current through
resistor \code{rA} in amperes, which corresponds to the current
$i_A$ in our problem.

Enter \code{ijb} in the entry line and press \key{ENTER}, then
$-5.4$ is shown in the history display, which is the current through
current source \code{jB} in amperes, corresponding to the current
$i_B$ in our problem.

Enter \code{ircc} in the entry line and press \key{ENTER}, then 6.\
is shown in the history display, which is the current through
resistor \code{rCC} in amperes, which corresponds to the current
$i_C$ in our problem.

This example illustrates why it is important that Symbulator can
find the currents through current sources: in this problem, the
current of source \code{jB} was unknown and required as an answer.
Let us see another example.

\problem{004}

\emph{Problem: For the circuit with two nodes shown in the figure,
find the currents $i_1$, $i_2$, $i_3$ and $i_4$.}

Solution:

As always, the first step is to name the circuit nodes. This circuit
has two nodes, or at least that is what the problem statement says.
Let us carefully examine the circuit. Clearly, we see that there are
six junctions which may be nodes. We also notice that the four outer
connections have no resistive elements or current sources: they are
short circuits. They have identical voltage at both ends. This means
that the six junctions are two groups of three, where each group has
the same voltage. When the problem says ``the circuit of two
nodes'', it refers to this fact. In terms of circuit voltages, these
six junctions are only two nodes. However, the problem also asks us
to find the current through the short circuits. Current can enter a
node and also leave it, but it cannot flow \emph{along} the node.
So, for the purposes of simulation, we must consider those
connections as short circuits, not nodes. This requirement is
identical for any simulator, including SPICE and Symbulator. So,
this circuit has six nodes and four short circuits, as well as other
elements.

We will also use this example to show how to define a
voltage-dependent current source whose control voltage is the
voltage between two nodes. Which of these six nodes will be the
reference node? Any one of them is a possibility, but we prefer to
choose as a reference node the one which is between the dependent
source and the 2.5 ohm resistor. Why? Because if we use that node as
the reference, then the control voltage of the dependent source is
simplified, for it can be defined as a single node voltage. Suppose
that we name the nodes left to right as follows: 1 is the first
node, 0 is the second, 2 is the third, 3 is the fourth, 4 is the
fifth and 5 is the sixth node. With these node names, the control
voltage of the dependent source is \code{(v1-v0)}, and since
\code{v0} $= 0$, the control voltage is \code{v1}.

The second step is to enter the simulation command and circuit
description in the entry line.

How shall we choose the best directions for the current flow? The
drawing shows the direction for the current sources. For the
resistors, any direction will work. For the short circuits, remember
what we learned in the previous problem: it is best to make the
current flow direction agree with the directions of the problem
questions and answers. So, we define the current flow directions
according to the directions of the unknown currents.

In order to answer the questions more easily, we will name the short
circuits \code{s1}, \code{s2}, \code{s3} and \code{s4}, in the same
order as the currents $i_1$, $i_2$, $i_3$ and $i_4$.

This, then, is the simulation command and the circuit description:

\begin{entry}
sq\dc("r1,1,0,25; j1,0,2,.2*v1; r2,2,3,10; j2,4,3,2.5;
      r3,4,5,100; s1,1,2,0; s2,2,4,0; s3,0,3,0; s4,3,5,0")
\end{entry}

Two things are worth noticing about the circuit description. First,
we needed to place zeros at the end of each short circuit element
description, so that all the rows have the same length. Second, note
that the control voltage for the source is called \code{v1}, just as
in the drawing. This is because we intentionally assigned the
reference node at the appropriate place, and we named the other
terminal 1. Alternatively, had we named the two nodes \code{m} and
\code{n}, the control voltage would have been \code{(vm-vn)}.
Another alternative, just as easy, would have been to define the
control voltage as \code{vr1}, that is, the voltage drop across
\code{r1}, because \code{r1} is connected to nodes 1 and 0, in the
right sequence.

Start the simulation by pressing \key{ENTER} after you have entered
the circuit description. After the status messages are shown, the
history display shows `Done'. My calculator ran this simulation in
32 seconds. The time has increased over that of previous examples,
because there are more nodes. Since Symbulator uses nodal analysis,
the number of equations and unknowns increases with the number of
nodes. The simulation results are stored in the current folder.

Let us display the currents, with the same procedure we have already
seen many times. The current \code{is1} is $-2$.\ amperes,
\code{is2} is 3.\ amperes, \code{is3} is $-8$.\ amperes, and
\code{is4} is $-.5$ amperes.

Now let us look at another type of dependent source.

\subsection{Current-dependent current sources}

For current-dependent current sources, the control current is often
the current through an element. For example, if the element is
\code{r1}, then the current is expressed as \code{ir1}.

In some circuits, the control current is the current through a
branch with no elements. In this case, we need to define that branch
as a short circuit so we can get the current, just as in the
previous example. If the short circuit is called \code{s1}, then the
current will be \code{is1}.

The following problem shows an example of current-dependent current
sources.

\problem{005}

\emph{Problem: Given the following circuit, determine the value of
the voltage $v$ and the power dissipation of the independent current
source.}

Solution:

The first step is to name the circuit nodes and elements. We use the
bottom node as the reference node, so that voltage $v$ is equal to
the upper node voltage. If we had used the upper node as the
reference, then voltage $v$ would be the negative of the lower node
voltage. We name the 2~k$\Omega$ resistor \code{rx}, and define its
current as flowing upwards, so that the direction of the current in
\code{rx} matches that of the control current $i_x$. This makes it
easier to interpret the simulation results. The current flow
direction through the other 6~k$\Omega$ resistor is not important.

The second step is to enter the simulation command and circuit
description in the entry line:

\begin{entry}
sq\dc("r1,1,0,6E3; j1,0,1,2*irx; j2,0,1,24E-3; rx,0,1,2E3")
\end{entry}

Note that the control current is labeled $i_X$ in the drawing, but
it is called \code{irx} in the circuit description. Why did we use a
different name? The reason is very simple: Symbulator cannot know
what we mean by $i_X$, because it cannot see the circuit schematic!
All control voltages or currents must be in terms of the circuit
elements. Since \code{rx} is a circuit element, Symbulator knows
that \code{irx} is the current through that resistor.

In the value \code{6E3}, `E' indicates the power of 10, so
\code{6E3} means $6 \times 10^{3}$. This nomenclature is used in
most programmable calculators, including Texas Instruments, HP and
Casio.

Note that even though the circuit resistances are given in
k$\Omega$, the circuit description values must be in ohms. Currents
in the circuit are specified in mA, but we must specify the values
in A for the circuit description. Symbulator does not understand
unit prefixes such as k or m.

You must keep in mind that we must describe the circuit in terms
that Symbulator understands. This avoids absurd situations such as
defining the control current as $i_X$ or specifying values in mA,
which would result in the wrong answers.

Start the simulation by pressing \key{ENTER}. `Done' is shown in the
history display after the simulation status messages are shown. My
calculator took 10 seconds to simulate this circuit. The simulation
results are stored in the current folder. This example clearly
demonstrates the agility and potential of Symbulator as a tool for
learning and checking answers. It would take several minutes to
solve this circuit by hand and find all of the voltages, currents
and power dissipations, and there is a chance that we may make a
mistake. With Symbulator, obtaining all these results takes only ten
seconds, and we can be confident that the results are correct.

The answers are shown with the same process used in the previous
examples. Voltage $v$ in the schematic is equivalent to \code{v1},
which is 14.4~V. The power dissipation of the independent source is
called \code{pj2}, which is $-.3456$~W. Since the problem asked for
the power dissipated by the source, \code{pj2} is the answer. If the
problem had asked for the power generated by the source, we would
change the sign of the answer, and the generated power would be
.3456~W (positive). Let us do another example, to reinforce the
concepts learned so far.

\problem{006}

\emph{Problem: Determine $i_X$ in the following circuit.}

Solution:

As always, the first step is to name all the circuit nodes and
elements. Like Problem 004, this circuit seems to have two nodes,
when in fact we need more nodes and some short circuits to simulate
it. For this problem, the location of control current $i_1$ forces
us to view the upper part of the circuit as two nodes connected by a
short circuit, which we will call \code{s1}. On the other hand, the
location of the unknown current $i_X$ forces us to view the lower
part of the circuit as two more nodes, and to use another short
circuit which we will call \code{sxx}. We cannot use the name
\code{sx}, because that is a reserved system variable name.

With this circuit, it is irrelevant which node we decide to call the
reference node. The current flow direction in the resistors can also
be assigned in either direction. The current flow in the current
sources must, as always, be assigned in the direction shown in the
drawing. The current flow directions in the short circuits are
chosen so that they agree with the control current and the unknown
current.

The second step is to enter the simulation command and circuit
description in the entry line:

\begin{entry}
sq\dc("r1,1,0,5E3; j1,1,0,4E-3; s1,1,2,0; sxx,3,0,0;
      j2,2,3,3*is1; r2,2,3,20E3")
\end{entry}

It should not be a surprise at this point that the control current
is called $i_1$ in the drawing, but is called \code{is1} in the
description.

Start the simulation by pressing \key{ENTER}. `Done' is shown in the
history display when the simulation is finished. My calculator took
17.5 seconds to simulate this circuit. The results are stored in the
current folder.

The answer to the problem, the unknown current $i_X$, is obtained by
recalling \code{isxx}, which is .000571429 amperes.

\section{Voltage source}

Ideally, a voltage source is a circuit element which forces a
defined voltage difference between its two terminals, and generates
or dissipates power.

\subsection{Notation}

The minimum information needed to define a voltage source is:

\emph{What kind of element is it?} It is a voltage source. The
letter that identifies voltage sources is \code{e}.

\emph{How is the source named?} To distinguish one voltage source
from another, we must give each one a different name, and those
names all start with \code{e}.

\emph{Where is the source connected?} A voltage source has two ends,
each connected to a different node. It is necessary to specify both
nodes to which the source is connected, and the order in which we
specify the nodes also specifies the polarity of the source: the
first node with respect to the second. In the circuit drawing, the
polarity of the source is specified by the `plus' and `minus' signs
next to the source.

When entering a voltage source in the circuit description, the order
in which the nodes are listed defines the source polarity. However,
the polarity also indirectly specifies the direction of current flow
through the source, from the first node to the second node. The
voltage drop across the voltage source is measured from the first
node with respect to the second node.

This point may cause some confusion if you are not familiar with
circuit simulation. You may think that, if the first node of the
voltage source is positive with respect to the second node, then the
current must flow from the second node to the first. Indeed, it is
like that, in most cases.

The important point to understand is that the polarity used to
define a variable does not necessarily establish which node is at
the higher voltage. As an example, consider the following: ``The
voltage of node A is $-5$~V with respect to node B''. In this
example, the polarity is defined at node A with respect to node B,
but node B is at a higher voltage than node A.

Similarly, defining the direction of current flow through an element
does not set the actual direction of current flow. Consider this
phrase: ``the current flowing from node A to node B is $-3$~A''. In
this case, the current is defined as flowing from A to B, but the
real direction of conventional current flow is from B to A.

Summarizing: in all the elements we have so far used, the current
flow direction used to define the simulation results is from the
first node to the second node, and the polarity of the voltage drop
across an element is defined as the voltage of the first node with
respect to the second node. In addition, the voltage source value is
given as the voltage of the first node with respect to the second.

\emph{What is the value of the source?} The value must be specified
in volts. For a direct current simulation, the value of an
independent voltage source can be a real number or a variable. If
the source is dependent, its value must be an algebraic expression.
This expression will be composed of a number and a variable. This
variable will be one of the simulation answers.

So, in a circuit description, a voltage source is defined as:

\begin{entry}
ename, node 1, node 2, value in volts
\end{entry}

As we have seen, the value of the voltage source is defined as the
voltage at node 1 with respect to node 2. For example, if the value
is 10~V, then the simulator assumes that the voltage difference from
node 1 with respect to node 2 is 10 volts.

Note that the voltage source description has four terms. If other
circuit element descriptions have five terms, then we must add a
zero to the source description, for example:

\begin{entry}
ename, node 1, node 2, value, 0
\end{entry}

Take care that all values are entered in volts, and not some other
unit such as millivolts (mV), otherwise the results will be wrong.

For analyses other than direct current, the value may be a phasor, a
function of frequency or a function of time.

\subsection{Related answers}

There are three simulation results for a voltage source: the current
through the source, the voltage drop across the source, and the real
power consumed by the source. The current is saved in a variable
called \code{iname}, where \code{name} is the name of the source.
The voltage drop is stored in a variable called \code{vname}. The
real consumed power is stored in a variable called \code{pname}.

The current, as previously explained in detail, is defined as
flowing from the first node to the second node in the description.

The voltage drop is defined as the voltage of the first node with
respect to the second node in the description.

The power is returned as the consumed power. If the source generates
power, the returned result is negative. You must remember this when
you interpret the simulation results.

\problem{007}

\emph{Problem: In the following circuit, determine the power
consumed by each of the elements.}

Solution:

The first step is to name all the nodes and elements. We will use
the bottom node as the reference node. Remember that voltage source
names must begin with \code{e} and the order of the nodes must
correspond to the polarity shown in the drawing. The order of the
nodes of the resistors is irrelevant.

The second step is to enter the simulation command and circuit
description in the entry line:

\begin{entry}
sq\dc("e1,1,0,120; r1,1,2,30; e2,2,3,30; r2,3,0,15")
\end{entry}

Start the simulation by pressing \key{ENTER}. After the status
messages are shown, `Done' is shown in the history display. My
calculator took 13 seconds to simulate this circuit. The answers are
stored in the current folder.

The power consumed by source \code{e1} is in the variable
\code{pe1}, and is $-240$ watts. This means that the source
generated 240~W, since the value is negative. However, since the
problem asked for the power consumption of each element, the answer
is $-240$~W. The power consumed by source \code{e2} is in variable
\code{pe2}, and is 60~W. Since the value is positive, the source
consumes power; it does not generate it. The power dissipated by
\code{r1} is in variable \code{pr1}, and is 120~W. The power
dissipated by \code{r2}, from variable \code{pr2}, is 60~W. With
this example, we can verify that Symbulator complies with the energy
conservation law. If we enter \code{pe1+pr1+pe2+pr2}, to find the
net power dissipation, the result is 0~W, as expected.

\subsection{Dependent sources}

All the concepts that we learned for the simulation of dependent
current sources apply to dependent voltage sources as well. Let us
see some examples for demonstration. We will start with a simple
problem with a single voltage-dependent voltage source.

\problem{008}

\emph{Problem: In the following circuit, determine the power
absorbed by each element.}

Solution:

The first step is to name all the nodes and elements. The order of
the nodes in the voltage sources is determined by the polarity shown
in the drawing. With the experience we now have, it is obvious that
if we use the bottom node as the reference node, and call the upper
left node \code{x}, then we can call the control voltage \code{vx}.
The other nodes are named 1, 2 and 3, moving clockwise around the
drawing.

The second step is to enter the simulation command and circuit
description in the entry line:

\begin{entry}
sq\dc("r1,x,0,30; e1,1,x,12; r2,1,2,8; r3,2,3,7; e2,3,0,4*vx")
\end{entry}

Start the simulation by pressing \key{ENTER}. After the status
messages are shown, the word `Done' is shown in the history display.
My calculator took 18 seconds to simulate this circuit. The results
are stored in the current folder.

If we enter \code{pr1} in the entry line and press \key{ENTER}, the
power dissipated by resistor \code{r1} is shown as 96/125~W. As
explained, the simulation results are exact, because all of the
circuit values are exact. Usually, textbook answers are not given as
exact values like 96/125~W, but as approximate values like 0.768~W.
The exact value is as correct as the approximate value;
nevertheless, in an examination the professor may prefer the
approximate value. Note that the terms `exact' and `approximate' are
used in the same sense as defined by the operation of programmable
calculators.

It is simple to obtain the approximate answer instead of the exact
answer: just press \key{$\diamond$} before pressing \key{ENTER}.

With \code{pr1} in the entry line, if we press \key{$\diamond$}
\key{ENTER}, then the approximate result of .768~W is shown. This
same procedure is used for the other variables.

Enter \code{pe1}, then press \key{$\diamond$} \key{ENTER} and the
result of 1.92~W is shown.

Enter \code{pr2}, then press \key{$\diamond$} \key{ENTER} and the
result of .2048~W is shown.

Enter \code{pr3}, then press \key{$\diamond$} \key{ENTER} and the
result of .1792~W is shown.

Enter \code{pe2}, then press \key{$\diamond$} \key{ENTER} and the
result of $-3.072$~W is shown.

Next, we will do an example with two voltage-dependent voltage
sources, with more complicated control voltages.

\problem{009}

\emph{Problem: Find the power absorbed by each element of this
circuit.}

Solution:

The first step is to name all the circuit nodes and elements. The
order of the nodes in the circuit descriptions for the three voltage
sources is determined by the polarities shown in the drawing.
Resistor polarities are irrelevant. There is no clear choice for the
reference node, so we choose the bottom left node as the reference,
and name the other nodes in clockwise order as \code{a}, \code{b},
\code{c}, \code{d} and \code{e}.

The second step is to enter the simulation command and the circuit
description:

\begin{entry}
sq\dc("e1,a,0,40; r1,a,b,5; r2,b,c,25; r3,c,d,20;
      e2,e,d,2*vr3+vr2; e3,e,0,4*vr1-vr2")
\end{entry}

Start the simulation by pressing \key{ENTER}. After the status
messages are shown, `Done' is shown in the history display. My
calculator took 26 seconds to simulate this circuit. The answers are
stored in the current folder.

For \code{pe1}, we obtain 80 watts.

For \code{pr1}, we obtain 20 watts.

For \code{pr2}, we obtain 100 watts.

For \code{pr3}, we obtain 80 watts.

For \code{pe2}, we obtain $-260$ watts.

For \code{pe3}, we obtain $-20$ watts.

These are the correct answers. Now let us see an example of a
dependent voltage source controlled by a circuit current.

\problem{010}

\emph{Problem: Find the current $i_X$ in this circuit.}

Solution:

The first step is to name the circuit nodes and elements. The order
of the nodes in the three sources is determined by the drawing. The
current flow direction through the 2 ohm resistor (which will be
named \code{rx}) must agree with that of $i_X$; the current flow
direction through the other resistor is irrelevant. We choose the
bottom node as the reference node, and name the other nodes 1, 2 and
3, moving clockwise around the circuit.

The second step is to enter the simulation command and the circuit
description:

\begin{entry}
sq\dc("e1,1,0,10; rx,1,2,2; j1,0,2,3; r1,2,3,1; e2,3,0,2*irx")
\end{entry}

Start the simulation by pressing \key{ENTER}. After the status
messages are shown, `Done' appears in the history display. My
calculator took 16 seconds to simulate this circuit. The answers are
stored in the current folder.

Enter \code{irx}, then press \key{$\diamond$} \key{ENTER} to obtain
the current of 1.4~A through \code{rx}, which is correct. We will
finish this chapter with a problem that combines everything we have
learned so far.

\problem{011}

\emph{Problem: Find the currents, voltage drops and power
dissipation for each element in this circuit. Show that the sum of
the powers is zero.}

Solution:

The first step is to name all the circuit nodes and elements. By
this point, you should be able to do this without additional
explanation. The second step is to enter the simulation command and
the circuit description:

\begin{entry}
sq\dc("e1,1,0,60; r1,1,0,20; e2,1,2,5*ir1; j1,2,0,ve1/4; r2,2,0,5")
\end{entry}

Start the simulation by pressing \key{ENTER}. `Done' is shown in the
history display after the status messages are shown. My calculator
took 15 seconds to simulate this circuit. The answers are stored in
the current folder. The results are: \code{ve1} and \code{vr1} are
both 60 volts, \code{ve2} is 15 volts, \code{vj1} and \code{vr2} are
both 45 volts, \code{ie1} is $-27$ amperes, \code{ir1} is 3 amperes,
\code{ie2} is 24 amperes, \code{ij1} is 15 amperes, \code{ir2} is 9
amperes, \code{pe1} is $-1620$ watts, \code{pr1} is 180 watts,
\code{pe2} is 360 watts, \code{pj1} is 675 watts, \code{pr2} is 405
watts. Note that only one source is generating power. How long would
it take a student to find, without errors, these 15 answers? It took
the Symbulator, on average, 1 second for each one.

\section{Correcting a wrong description}

Suppose that we make an error in entering the circuit description,
but we do not notice it until the simulation is complete. We need to
correct the error and run the simulation again. First, we must erase
the result variables in the current folder. Next, we must put the
simulation command and the circuit description in the entry line,
again. It is a tedious task to completely re-enter the circuit
description a second time. The skilled user will take advantage of
the fact that the circuit description is already written. There are
two ways to do this:

1) If the previous description is still in the entry line, it will
be shown as white characters on a dark background. Move the cursor
to the error with the blue cursor keys and correct it. Note that the
\key{BACKSPACE} key is used to erase a character to the left of the
cursor.

2) If, for some reason, the circuit description is not in the entry
line, move the cursor upwards until the circuit description in the
history area is highlighted. Press \key{ENTER} to make a copy in the
entry line. The error in this copy is corrected as described above.

Once the error is corrected, run the simulation again by pressing
\key{ENTER}. With this explanation, we have finished the chapter on
numerical direct current simulation.
